\documentclass[11pt,a4paper]{article}

\usepackage[utf8]{inputenc}
\usepackage[T1]{fontenc}
\usepackage{lmodern}
\usepackage[margin=1in]{geometry}
\usepackage{booktabs}
\usepackage{array}
\usepackage{tabularx}
\usepackage{longtable}
\usepackage{listings}
\usepackage{xcolor}
\usepackage{hyperref}
\usepackage{amsmath,amssymb}
\usepackage{enumitem}
\usepackage{float}
\usepackage{caption}

\hypersetup{
    colorlinks=true,
    linkcolor=blue!60!black,
    citecolor=blue!60!black,
    urlcolor=blue!60!black,
}

\lstdefinelanguage{tlaplus}{
    morekeywords={LET,IN,IF,THEN,ELSE,EXCEPT,RECURSIVE,UNCHANGED,CHOOSE,
                  VARIABLES,CONSTANTS,EXTENDS,MODULE,INIT,NEXT,TRUE,FALSE,
                  INVARIANT,INVARIANTS,CONSTANT},
    sensitive=true,
    morecomment=[l]{\\*},
    morestring=[b]",
    literate={|->}{{$\mapsto$}}2
             {\\in}{{$\in$}}1
             {\\A}{{$\forall$}}1
             {\\E}{{$\exists$}}1
             {/\\}{{$\land$}}1
             {\\/}{{$\lor$}}1
             {=>}{{$\Rightarrow$}}1
             {->}{{$\rightarrow$}}1
             {<<}{{$\langle$}}1
             {>>}{{$\rangle$}}1
             {!=}{{$\neq$}}1
             {<=}{{$\leq$}}1
             {>=}{{$\geq$}}1,
}

\lstset{
    basicstyle=\ttfamily\small,
    keywordstyle=\bfseries\color{blue!70!black},
    commentstyle=\itshape\color{gray},
    stringstyle=\color{red!70!black},
    breaklines=true,
    frame=single,
    framerule=0.5pt,
    rulecolor=\color{gray!50},
    xleftmargin=1em,
    framexleftmargin=0.5em,
    numbers=none,
    aboveskip=0.8em,
    belowskip=0.8em,
}

\newcolumntype{L}[1]{>{\raggedright\arraybackslash}p{#1}}
\newcolumntype{C}[1]{>{\centering\arraybackslash}p{#1}}

% Allow more aggressive interword stretching to avoid overfull hboxes
% caused by long \texttt{...} identifiers that cannot be hyphenated.
\emergencystretch=3em

% Permit line breaks inside \texttt after "/" and "_" so that long
% file paths (e.g. matcher_lean/TheoremsElegant.lean) can wrap.
\makeatletter
\def\BreakableChars#1{\discretionary{}{}{}\ttfamily#1}
\makeatother

\title{From Bounded Model Checking to Mechanized Proof:\\
A Multi-Level Verification Case Study for a\\
Price-Time Priority Matching Engine\thanks{Artifacts available at
\url{https://github.com/formally-verified-exchange/verified-matching-engine}.}}

\author{%
  Ara Aslyan\\
  Strike Technologies, New York\\
  \texttt{aaslyan@striketechnologies.com}
}
\date{04/11/2026}

\begin{document}

\maketitle

\begin{abstract}
We present a case study in multi-level formal verification of a
feature-rich price-time priority matching engine---the software at the
heart of every stock, futures, and cryptocurrency exchange. The engine
supports five order types (\textsc{limit}, \textsc{market},
\textsc{market-to-limit}, \textsc{stop\_limit}, \textsc{stop\_market}),
iceberg display quantities, four self-trade prevention policies,
post-only, fill-or-kill, and minimum-quantity semantics.

The paper is structured around the assurance layers and the defects they
made visible. Bounded model checking with TLC, across configurations
totaling 36 million reachable states, produced a concrete FIFO
counterexample: triggered stop orders retained their original submission
timestamp. The same modeling pass exposed a prose-specification gap in
the interaction between STP DECREMENT and iceberg reload semantics.
Mechanized proof found two unbounded/transcription issues: the matching
loop's original constant fuel bound of 100 steps made the key
preservation theorem false for books with more than 100 contra-side
orders, and the full-invariant extension exposed an MTL/minimum-quantity
routing error in the Lean and TLA+ transcriptions. Finally, replaying
TLC-generated traces against a C++ implementation found a dispose defect:
a triggered stop-market order with a persistent time-in-force could
leave a market remainder resting on the book.

We prove in Lean~4 that the processing pipeline preserves the full
book-state invariant suite of the specification's \S13---not only that
the book stays uncrossed and sorted on both sides (\texttt{AllInv}), but
also no ghost orders, status consistency, FIFO within each price level,
and no resting market, market-to-limit, or minimum-quantity orders,
together with the post-only and self-trade-prevention guarantees on
emitted trades---for arbitrary book sizes, with a fuel bound computed
from the book state and proved sufficient. The proof spans three files
totaling roughly 5{,}700 lines: a 3{,}965-line constructive core, a
1{,}423-line extension lifting the remaining \S13 invariants, and a
complementary 311-line structural proof via three-case analysis of the
matching termination condition.

The overarching lesson is that feature composition is where bugs hide.
Each individual feature is correctly specified in isolation. The bugs
live in intersections that prose review cannot see but an operational
model exposes by construction.
\end{abstract}


% ============================================================================
\section{What Is a Matching Engine?}
\label{sec:intro}
% ============================================================================

\subsection{The Role of a Matching Engine}

When you buy a stock, someone else is selling it. The matching engine
finds the seller. Every exchange---stock market, futures exchange,
cryptocurrency platform---runs a matching engine at its core. Every
second, it processes thousands of incoming orders and decides which ones
trade with which, at what price, in what quantity, and in what sequence.

If the matching engine makes a mistake, real money moves wrong. An order
that should trade does not. An order that should not trade does. An
order sitting at the front of a queue is pushed behind a newer one.
These are not theoretical concerns: every such error is a financial
event with an identifiable winner and loser.

This paper is about formally verifying that a matching engine does what
it claims to do---not just for small test cases, but for all possible
inputs of all sizes.

\subsection{Orders}
\label{sec:orders-intro}

An order is an instruction to trade. The simplest example: ``I want to
buy 100 shares at \$150 or less.'' This is a \textsc{buy limit} order.
The engine will fill it if a willing seller exists at that price or
better; otherwise, it waits on the book.

The engine studied in this paper supports five order types:

\begin{itemize}[nosep]
\item \textbf{\textsc{limit}}: rest at a specified price; fill at that
      price or better. ``Buy 100 at \$150.'' Will wait indefinitely for
      a matching seller.

\item \textbf{\textsc{market}}: fill at whatever price is available
      right now; never rest on the book. ``Buy 100 at any price.''
      Executes immediately or not at all.

\item \textbf{\textsc{market-to-limit} (MTL)}: aggressive first, then
      patient. Match at any available price; lock to the first execution
      price; rest any unfilled remainder as a limit order at that price.

\item \textbf{\textsc{stop\_limit}}: dormant until a trigger event,
      then active. ``Once the last trade price reaches \$155, enter a
      buy limit order at \$156.'' Used for stop-loss and momentum
      strategies.

\item \textbf{\textsc{stop\_market}}: same trigger logic, but converts
      to a market order on activation rather than a limit order.
\end{itemize}

Orders also carry \emph{modifiers} that change their behavior:

\begin{itemize}[nosep]
\item \textbf{Iceberg} (display quantity): only a small ``visible
      slice'' is shown to the market. When that slice fills, the order
      reloads a fresh slice from its hidden reserve---but loses its
      place in the queue and joins the back of the line at its price.

\item \textbf{Post-only}: the order must not aggress against the
      opposite side. If it would immediately trade, it is rejected.
      Useful for market makers who want to add liquidity, not consume
      it.

\item \textbf{Fill-or-Kill (FOK)}: fill the entire order immediately,
      or cancel it entirely. No partial fills, no resting.

\item \textbf{Minimum quantity (MinQty)}: the first fill must be at
      least this large; after the initial threshold is met, the
      remainder may rest and fill in smaller pieces.

\item \textbf{Self-trade prevention (STP)}: four policies govern what
      happens when an incoming order would trade against a resting order
      from the same participant. Options: cancel the incoming order
      (\textsc{cancel\_newest}), cancel the resting order
      (\textsc{cancel\_oldest}), cancel both
      (\textsc{cancel\_both}), or reduce both quantities by the
      overlapping amount (\textsc{decrement}).
\end{itemize}

\subsection{The Order Book}
\label{sec:book-intro}

The order book is the matching engine's central data structure. It has
two sides: bids (buy orders) and asks (sell orders). Each side is
organized by price, with the most competitive prices at the top.

\begin{center}
\begin{lstlisting}[language={},basicstyle=\ttfamily\footnotesize,frame=single,xleftmargin=0pt,framexleftmargin=0pt]
         BIDS (buyers)          ASKS (sellers)
  Price    Qty   Orders    |  Price    Qty   Orders
  ------  -----  -------   |  ------  -----  -------
  $151     200   [B, C]    |  $152     100   [D]
  $150     500   [A]       |  $153     300   [E, F]
  $149     150   [G]       |  $155     200   [H]

         Best bid: $151          Best ask: $152
         Spread: $1
\end{lstlisting}
\end{center}

In this book, the best bid is \$151 and the best ask is \$152. The
\emph{spread} is \$1. No immediate trade is possible because no buyer
is willing to pay as much as the cheapest seller is asking.

When a new buy order arrives at \$152 or higher, it \emph{crosses the
spread}: it is willing to pay the ask price, so a trade occurs. The
matching engine fills the incoming order against the cheapest resting
ask, consuming quantity from the front of that level's queue, and
repeating until the incoming order is satisfied or no crossable orders
remain.

Within each price level, orders are served in \emph{first-in,
first-out} (FIFO) order: the order that arrived earliest has
highest priority. This is the ``time priority'' in price-time priority.
It means that among all sellers willing to sell at \$152, the one who
submitted first gets filled first.

\subsection{What Must Be True---Always}
\label{sec:invariants-intro}

A matching engine makes a set of guarantees. These must hold after
every single operation, no exceptions:

\begin{enumerate}[nosep]
\item \textbf{The book is not crossed.} The best bid is always strictly
      less than the best ask (or one side is empty). If the bid were
      higher than the ask, a trade should have occurred.

\item \textbf{The book is sorted.} Bids are arranged in strictly
      descending price order; asks in strictly ascending order. A buyer
      willing to pay more always has higher priority than one willing to
      pay less.

\item \textbf{FIFO within each price level.} Among orders at the same
      price, earlier arrivals are served first. Timestamps are
      strictly increasing down each level's queue.

\item \textbf{No ghost orders.} Every resting order has a positive
      remaining quantity. An order with zero quantity left must have
      been removed.

\item \textbf{Post-only safety.} A post-only order never results in a
      trade. If it would have aggressed, it was rejected before any
      trade occurred.

\item \textbf{Self-trade prevention.} No trade occurs between two
      orders belonging to the same STP group.
\end{enumerate}

These rules sound obvious. They are obvious---individually. The
difficulty is maintaining all of them simultaneously when multiple
features interact. Guaranteeing rule~3 (FIFO) requires careful handling
of every path by which an order can enter the book. Guaranteeing
rule~6 (STP) requires careful handling of every order type, including
stop orders that re-enter the pipeline from a dormant state. Features
that interact with each other create combinatorial cases that no single
human reviewer is likely to hold in mind simultaneously.

\subsection{Where Bugs Hide: Feature Interactions}
\label{sec:feature-interactions}

Each feature in a matching engine can be specified correctly in
isolation. The danger is interactions. Consider:

\begin{itemize}[nosep]
\item Stop orders are correctly described: they are dormant until
      triggered, then converted and processed.
\item FIFO insertion is correctly described: orders are appended to the
      back of their price level's queue, with a timestamp reflecting
      their arrival time.
\end{itemize}

These two descriptions are each locally correct. But they interact:
when a stop order is triggered, it is converted and re-enters the
pipeline. If its timestamp is not updated at trigger time, it carries
its original submission timestamp---which may be earlier than orders
already resting on the book. It will be appended to the back of the
queue, but its timestamp will be out of order. FIFO is violated.

A human reviewer reading the stop section and the FIFO section
separately will not see this. Both sections are individually correct.
The interaction is invisible unless you hold both simultaneously.

A second example: iceberg reload is correctly described for the fill
path. But STP DECREMENT reduces a resting order's visible quantity
without generating a trade---so it is not the fill path. If DECREMENT
drives visible quantity to zero while hidden quantity remains, the
iceberg never reloads, and the order is permanently stranded: it sits
on the book, cannot match, and cannot reload.

This is the problem we set out to solve. The goal is a verification
approach that does not read features section by section, but evaluates
all of them simultaneously. The rest of this paper describes what we
built and what we found.


% ============================================================================
\section{How We Checked It---Overview}
\label{sec:overview}
% ============================================================================

We verified the engine at two complementary levels, each with a
different scope and a different kind of guarantee.

\subsection{Level 1: Try Every Combination (Bounded Model Checking)}
\label{sec:overview-tlc}

The first approach is exhaustive state-space exploration. We encoded
the engine's rules as a mathematical model in TLA+ and used the TLC
model checker to enumerate \emph{every possible sequence of events}
within a set of small configurations: up to 2 orders, quantities up to
2, prices drawn from a set of 2 or 3 values.

TLC is a brute-force tool. It generates every reachable state of the
model and checks that each one satisfies every stated invariant. If any
invariant is ever violated, TLC returns a concrete counterexample trace
showing exactly which sequence of orders produced the violation.

Within the configurations we explored, TLC generated 36 million
distinct states. \textbf{This produced one concrete invariant
counterexample and exposed one additional specification gap.}

The limitation of this approach is its scope. Production matching
engines hold millions of orders at prices across a continuous range. No
model checker can enumerate those states. A clean pass means ``no
violation found in the explored region,'' not ``no violation exists.''

\subsection{Level 2: Prove It for All Sizes (Theorem Proving)}
\label{sec:overview-lean}

The second approach is mechanized proof. We implemented the engine in
Lean~4---a programming language that also serves as a proof assistant---and
wrote a machine-checked proof that the engine's key guarantee holds not
for small configurations, but for \emph{all possible inputs of all
sizes}.

A Lean proof is verified by the proof checker on every build. There is
no ``trust me, I checked it carefully.'' Either the proof compiles or
it does not. When it compiles, the guarantee is unconditional within
the stated preconditions.

\textbf{This process discovered two further issues}: the matching loop's
original termination bound was an arbitrary constant, which made the
key theorem false for large books, and the later full-invariant proof
exposed an MTL/minimum-quantity routing error in the executable
transcriptions. The fuel issue required a universally quantified
statement---bounded model checking had never reached configurations
large enough for the constant to matter.

\subsection{What We Found}
\label{sec:findings-summary}

\begin{center}
\small
\begin{tabular}{@{}lL{4cm}L{7cm}@{}}
\toprule
Finding & Tool & Description \\
\midrule
1 & Bounded model checking (TLC) & FIFO violation: triggered stop orders retain their original submission timestamp, producing a non-monotonic queue when the triggered order rests \\
\addlinespace
2 & TLA+ modeling / specification review & Iceberg stranding: STP DECREMENT can drive visible quantity to zero without triggering iceberg reload, permanently stranding the order; the checked TLA+ model includes the fix, so this is reported as a prose-specification gap rather than a TLC counterexample \\
\addlinespace
3 & Theorem proving (Lean~4) & Unsound fuel bound: the matching loop's constant limit of 100 steps makes the preservation theorem false for books with more than 100 contra-side orders \\
\addlinespace
4 & Full-invariant proof (Lean~4) & Transcription defect: a passing MinQty pre-check could route a \textsc{market-to-limit} order directly to dispose before conversion, allowing it to rest as \textsc{mtl}; the prose specification was correct, but both executable transcriptions diverged \\
\addlinespace
5 & TLA+-trace conformance testing & C++ dispose defect: a triggered \textsc{stop\_market} order retaining \textsc{day}/\textsc{gtd} time-in-force could rest a market remainder on the book, violating \textsc{NoRestingMarkets} \\
\bottomrule
\end{tabular}
\end{center}

\subsection{What Each Layer Establishes}
\label{sec:assurance-boundary}

The three methods verify different objects and deliver different kinds
of guarantee. Because ``verified'' and ``proved'' are easy to over-read,
Table~\ref{tab:assurance} states precisely what is established, by which
artifact, and---equally important---what is not.

\begin{table}[H]
\centering
\caption{The assurance boundary. Each layer verifies a distinct object;
the guarantees do not compose into an end-to-end machine-checked
refinement from prose specification to C++.}
\label{tab:assurance}
\small
\begin{tabular}{@{}lL{3.9cm}L{6.4cm}@{}}
\toprule
Layer & Object verified & Guarantee \\
\midrule
TLC (TLA+)          & The TLA+ model & Exhaustive, but only within the finite configurations of Table~\ref{tab:configs} (up to 2--3 orders, 2--3 prices). No claim beyond those bounds. \\
\addlinespace
Lean~4              & The Lean executable reference implementation & The full \S13 book-state invariant suite is preserved for \emph{arbitrary} finite books, under preconditions discharged by well-formedness. Machine-checked on every build. \\
\addlinespace
Conformance testing & The C++ implementation & Agreement with the TLA+ model on projected observables, for the traces actually replayed (\S\ref{sec:conformance-pipeline}). Testing, not proof. \\
\addlinespace
\emph{Not established} & Lean reference $\leftrightarrow$ C++ engine & No general refinement proof links the proved Lean reference to the shipped C++ engine; the two are connected only by conformance testing. \\
\bottomrule
\end{tabular}
\end{table}

The Lean guarantee is unbounded but concerns the Lean reference; the C++
guarantee concerns the shipped engine but is bounded by the traces
replayed. Neither substitutes for the other, and no layer claims a
machine-checked path from the prose specification all the way to C++.


% ============================================================================
\section{What Model Checking Found}
\label{sec:findings}
% ============================================================================

\subsection{Bug 1: The Stale Timestamp}
\label{sec:bug1}

\paragraph{The scenario.} A stop order is submitted early---say at
time~2---with a stop price of \$1. It goes dormant, waiting for the
last trade price to reach \$1. Meanwhile, another order arrives at
time~3 and rests on the bid side at \$1.

Later, a trade occurs at \$1, triggering the stop order. The stop
converts to a live limit order and enters the matching pipeline. It
does not fill immediately (the ask side is now empty), so it rests on
the bid side at \$1. It is appended to the back of that price level's
queue---behind the order that arrived at time~3.

The queue at price \$1 now reads: [order at time~3, order at time~2].
But time~3 $>$ time~2, so the queue is not in ascending timestamp order.
FIFO is violated.

\paragraph{The exact trace.} TLC produced the following minimal
counterexample:

\begin{figure}[H]
\centering
\begin{lstlisting}[language={},basicstyle=\ttfamily\footnotesize,frame=single,xleftmargin=0pt,framexleftmargin=0pt]
Step 1: Submit SELL LIMIT  qty=1, price=1      (id=1, ts=1)
        -> Rests on ask:  askQ[1] = [order1(ts=1)]

Step 2: Submit BUY STOP_LIMIT  stop=1, price=1, qty=1  (id=2, ts=2)
        -> lastTradePrice = NULL -> does not trigger
        -> Placed in stops = {order2(ts=2)}

Step 3: Submit BUY LIMIT  qty=2, price=1       (id=3, ts=3)
        -> Matches askQ[1]: fills order1 fully (fillQty=1)
        -> Trade: price=1, aggressor=3, passive=1
        -> lastTradePrice = 1
        -> order3: remainingQty=1, rests on bid
           bidQ[1] = [order3(ts=3)]
        -> Stop cascade: lastTradePrice(1) >= stopPrice(1)
           -> order2 triggers, converts to BUY LIMIT
           -> Timestamp NOT updated: ts=2 (STALE)
           -> Ask side empty, no fill
           -> Dispose: order2 rests on bid
              bidQ[1] = [order3(ts=3), order2(ts=2)]

        FIFOWithinLevel check:
            bidQ[1][1].timestamp = 3
            bidQ[1][2].timestamp = 2
            3 < 2 ?  FALSE  ->  INVARIANT VIOLATED
\end{lstlisting}
\caption{TLC counterexample trace for the FIFO violation. The triggered
stop order (id=2) retains its original timestamp (ts=2) and is appended
behind an order with a later timestamp (ts=3), producing a
non-monotonic queue.}
\label{fig:trace}
\end{figure}

\paragraph{Why manual review missed this.} The prose specification
treats stop triggering (\S10.1) and order insertion (\S6.1) as separate
sections. Each section is locally correct. The interaction---that
converted stops carry stale timestamps into insertion---is invisible
when reading either section in isolation. TLC holds all three
simultaneously by construction.

\paragraph{The fix.} One line in the stop-trigger procedure: assign a
fresh timestamp before re-entering the pipeline.

\begin{lstlisting}[language=tlaplus]
converted == [ConvertStop(s) EXCEPT !.timestamp = tm]
\end{lstlisting}

In prose: after converting a stop to its base type and before calling
\texttt{PROCESS}, assign the stop a fresh timestamp reflecting the
current time. This ensures the triggered order's queue priority reflects
when it entered the live matching pipeline, not when it was originally
submitted.

\subsection{Bug 2: The Stranded Iceberg}
\label{sec:bug2}

\paragraph{The scenario.} A resting iceberg order has total quantity~2
and display quantity~1, so its visible slice is~1 and its hidden
reserve is~1. An incoming order from the same STP group arrives with
STP policy \textsc{decrement}.

DECREMENT reduces both sides by the minimum of the two visible
quantities, without generating a trade. It reduces the resting order's
visible quantity by~1 (to~0) and its remaining quantity by~1 (to~1).
The incoming order is consumed and cancelled.

The resting order now has visible quantity~0 and remaining quantity~1.
It should reload its visible slice from its hidden reserve. But the
reload logic is only specified in the fill path. DECREMENT is not the
fill path. The reload never triggers.

The result is a permanently stranded order: visible quantity~0 means
\texttt{fillQty = min(incoming, 0) = 0}, so matching never touches it.
It sits on the book indefinitely, occupying a queue slot, blocking
nothing, and unreachable except by explicit cancel.

\paragraph{The illustrative state.}

\begin{lstlisting}[language={},basicstyle=\ttfamily\footnotesize,frame=single,xleftmargin=0pt,framexleftmargin=0pt]
Resting iceberg: qty=2, displayQty=1, visibleQty=1,
                 remainingQty=2, stpGroup="G1"        (id=R)
Incoming:        qty=1, stpGroup="G1", stpPolicy=DECREMENT  (id=I)

DECREMENT:
    reduceQty = min(1, 1) = 1
    I.remainingQty = 0   -> I is cancelled
    R.visibleQty   = 0
    R.remainingQty = 1   -> R remains on book

Post-decrement state of R:
    visibleQty=0, remainingQty=1, displayQty=1
    -> Reload path not invoked
    -> R sits on book, invisible, unreachable by matching
\end{lstlisting}

\paragraph{Note on claim level.} This scenario is reachable in the
prose specification as written. However, the TLA+ model as checked
incorporates the fix because the stranded state would cause the model
checker to loop: subsequent orders would find visible quantity~0 and
skip the order without decrementing fuel, creating an effective infinite
loop in the matching operator. The bug is presented as a specification
gap---a path the spec silently omits---rather than as a TLC invariant
violation with a counterexample trace.

\paragraph{The fix.} After DECREMENT, check whether the resting order
is an iceberg with visible quantity~0 and hidden reserve remaining. If
so, reload:

\begin{lstlisting}[language=tlaplus]
ELSE IF newRest.visibleQty = 0 /\ newRest.displayQty /= NULL THEN
    LET reloadQty == Min(newRest.displayQty, newRest.remainingQty)
        reloaded  == [newRest EXCEPT !.visibleQty = reloadQty,
                                     !.timestamp = tm,
                                     !.status = "PARTIAL"]
        sc == SetContra(side, bQ, aQ, bestP,
                        restTail \o <<reloaded>>)
    IN DoMatch(newInc, sc.bQ, sc.aQ, tds, tm + 1, fuel - 1)
\end{lstlisting}

This mirrors the reload behavior in the normal fill path, extending it
consistently to the DECREMENT path.

\subsection{What Else We Noticed}
\label{sec:observations}

Three additional interactions were noted during model exploration as
potential specification ambiguities. No invariant violation was produced
for them; they are reported as design observations.

FOK combined with STP DECREMENT creates an ambiguity: the FOK pre-check
counts non-conflicting visible quantity to ensure full fill is possible,
but DECREMENT can consume the incoming order's quantity without
generating trades. Whether the result constitutes a FOK violation
depends on whether ``fully filled'' means ``fully traded'' or
``remaining quantity reaches zero''---the spec does not address this
case. Similarly, MinQty combined with STP DECREMENT may produce a
situation where the minimum-quantity pre-check passes but the actual
traded quantity falls below the minimum because DECREMENT consumed
quantity before any fills occurred. Finally, the spec clears
\texttt{minQty} in two separate places (the MTL path and the normal
matching path)---a duplication that is currently correct but fragile
under any future addition of routing paths.

\subsection{The Scope and Limits of Model Checking}
\label{sec:tlc-limits}

The configurations we explored generated 36 million distinct states.
This is large enough to find the feature-interaction bugs above. It is
not large enough to reason about production-scale books.

A production matching engine may hold thousands of resting orders per
price level. Model checking with 2 orders, quantities up to 2, and 3
prices cannot say anything about whether a matching loop terminates
correctly with 200 orders at one price. For that kind of guarantee, we
need a proof.


% ============================================================================
\section{The Lean Proof---Key Ideas}
\label{sec:lean}
% ============================================================================

TLC searches a bounded state space. Production engines hold millions of
orders, far beyond any configuration TLC can enumerate. To establish
correctness at arbitrary scale, we translated the fixed specification
into Lean~4 and proved the key preservation property mechanically.

\subsection{The Goal}
\label{sec:lean-goal}

We want to prove: after processing any valid order on any valid book,
the book is still uncrossed and sorted on both sides. For any number of
orders, any book size, forever.

The formal statement defines the combined invariant
\begin{align*}
\texttt{AllInv}(b) \;\equiv\;{}
  & \texttt{BookUncrossed}(b) \;\land\; {}\\
  & \texttt{bidsSortedDescB(b.bids) = true} \;\land\; {}\\
  & \texttt{asksSortedAscB(b.asks) = true},
\end{align*}
and the main theorem is:

\begin{lstlisting}[language={},basicstyle=\ttfamily\footnotesize]
theorem processOrder_preserves_AllInv
    (fuel : Nat) (o : Order) (b : BookState)
    (hpok : OrderProcOk o)
    (hstops : StopsNoPostOnly b)
    (h : AllInv b) :
    AllInv (processOrder fuel o b).book
\end{lstlisting}

The sortedness conjuncts are not cosmetic. \texttt{BookUncrossed} alone
is not sufficient for the induction to close: the induction hypothesis
needs to know that both sides of the book are monotone in price. This
was established empirically---the proof refused to close under the
weaker hypothesis, at which point the invariant was strengthened and the
proof proceeded. The proof refined the invariant.

\subsection{The First Surprise: The Theorem Was False}
\label{sec:lean-fuel}

Lean's termination checker requires every recursive function to be
provably total. The matching loop \texttt{doMatch} is recursive, so the
original implementation passed the same explicit \emph{fuel} argument into
the matching loop, using a constant bound:

\begin{lstlisting}[language={},basicstyle=\ttfamily\footnotesize]
def defaultFuel : Nat := 100
def doMatch (fuel : Nat) ... : MatchResult :=
  match fuel with
  | 0       => /* fuel exhausted: return as-is */
  | fuel'+1 => /* one step, recurse with fuel' */
\end{lstlisting}

Using this constant as the matching-loop bound hides an assumption:
matching always completes within 100 steps.
That assumption does not hold at arbitrary scale. Consider a book with
200 resting asks at price~\$50 and one incoming buy at price~\$100.
\texttt{doMatch} consumes the first 100 asks, exhausts its fuel, and
returns---with the incoming order still having remaining quantity and
100 asks still resting at price~\$50. The disposition phase then inserts
the buy at \$100 on the bid side. Now the best bid is \$100 and the
best ask is \$50. The book is crossed. \texttt{AllInv} is false.

This is precisely where bounded model checking and theorem proving
diverge. TLC's explored configurations never exceeded a handful of
orders per side, so the constant bound was never a binding constraint.
Only the universally quantified Lean statement forced the question: is
the fuel bound sufficient, or is it just an assumption?

\subsection{The Core Insight: Every Step Makes Progress}
\label{sec:lean-fuel-fix}

The fix is to compute the fuel bound from the book state and prove that
it is always sufficient:

\begin{lstlisting}[language={},basicstyle=\ttfamily\footnotesize]
def matchMeasure (contra : List PriceLevel) (_inc : Order) : Nat :=
    totalRemaining contra + orderCount contra + contra.length

def computeMatchFuel (b : BookState) (side : Side) : Nat :=
    matchMeasure (contraLevels b side) dummyInc + 1
\end{lstlisting}

The measure has three components. \texttt{totalRemaining}---the sum of
all remaining quantities on the contra side---accounts for fills.
\texttt{orderCount}---the number of individual orders---accounts for
orders that are skipped without a fill (STP cancellations, zero-visible
iceberg slots). \texttt{contra.length}---the number of price
levels---accounts for level exhaustion steps that remove an empty level.

The key technical lemma is the \emph{progress lemma}: every
non-terminal step of \texttt{doMatch} strictly decreases this measure.
Table~\ref{tab:branches} gives the branch-by-branch analysis.

\begin{table}[H]
\centering
\caption{Progress analysis for \texttt{doMatch} branches. Terminal
branches stop the loop; measure-decreasing branches recurse on a
strictly smaller state.}
\label{tab:branches}
\small
\begin{tabular}{@{}llL{5.5cm}@{}}
\toprule
Branch & Kind & Why the measure decreases \\
\midrule
Full fill (resting consumed)    & Decrease  & \texttt{orderCount} drops by 1 \\
Partial fill (incoming done)    & Terminal  & $\mathit{remainingQty}=0$ \\
STP cancel newest               & Terminal  & incoming cancelled \\
STP cancel oldest               & Decrease  & one order removed \\
STP cancel both                 & Terminal  & incoming cancelled \\
STP decrement                   & Decrease  & resting qty strictly reduced \\
Iceberg fill + reload           & Decrease  & fill reduced
                                              \texttt{totalRemaining} before reload \\
Level exhausted                 & Decrease  & empty level removed \\
Zero-visible skip               & Decrease  & one order removed \\
No crossable orders             & Terminal  & $\lnot\,\texttt{canMatchPrice}$ \\
\bottomrule
\end{tabular}
\end{table}

The iceberg reload case is the subtlest. Reload appears to ``add
liquidity back'': the resting order's visible quantity jumps from zero
to the display quantity. But \texttt{totalRemaining} is not affected by
this redistribution---only the visible/hidden split changes---and the
fill step that triggered the reload strictly decreased
\texttt{totalRemaining} by the filled quantity. The composition is a
net decrease.

\subsection{The Elegant Argument}
\label{sec:lean-elegant}

Once the constructive proof of the full pipeline was complete, a second
shorter proof was developed for the matching-and-dispose core, in the
file \texttt{TheoremsElegant.lean} (311 lines, 6 theorems).

The matching loop continues as long as the incoming order's price is
willing to cross the best price on the opposite side---a condition
captured by the predicate \texttt{canMatchPrice} (defined formally in
\S\ref{sec:invariants}). When the loop stops, this condition is false.

The insight is that when \texttt{doMatch} terminates on a
fuel-sufficient run, the resulting state falls into exactly one of
three cases:

\begin{enumerate}[nosep]
\item \textbf{Aggressor consumed.} Remaining quantity is zero (or the
      order was cancelled). The incoming order is gone.
      \texttt{doMatch} only removes orders from the contra side, never
      adds to it; removing orders cannot create a cross.

\item \textbf{Contra side empty.} Nothing left on the opposite side to
      cross against. An empty side is vacuously non-crossing.

\item \textbf{Non-crossing head.} Both sides are non-empty but the
      contra head does not cross the incoming price. This is the only
      way \texttt{doMatch} can stop with both sides active: if prices
      crossed, the loop guard \texttt{canMatchPrice} would fire and
      matching would take another step. Since it stopped, prices do not
      cross.
\end{enumerate}

There is no fourth case. A state where both sides are active and still
crossing is not a terminal state---fuel sufficiency rules it out. The
three cases compose with the disposition step directly: Case~1 returns
the book unchanged, Case~2 inserts on an empty contra side (vacuously
non-crossing), Case~3 inserts at a price strictly separated from the
contra head.

This 311-line proof treats \texttt{doMatch} as a black box
characterized by its termination condition. It reuses the constructive
proof's fuel-sufficiency lemma as an imported fact. The contrast
between the two proof styles is instructive:

\begin{center}
\small
\begin{tabular}{@{}lll@{}}
\toprule
                        & Constructive              & Structural              \\
\midrule
File                    & \texttt{Theorems.lean}    & \texttt{TheoremsElegant.lean} \\
Lines                   & 3{,}965                   & 311                     \\
Theorems/lemmas         & 74                        & 6                       \\
Argument                & Branch-by-branch          & Three-case analysis     \\
\bottomrule
\end{tabular}
\end{center}

The constructive proof discharges fuel sufficiency, sortedness
preservation, side isolation, and the mutual pipeline theorem. The
structural proof reuses those lemmas and shows that, once the
termination condition is established, the matching-and-dispose argument
itself is short.

\subsection{What Else the Proof Discovered}
\label{sec:lean-discovered}

Two things were discovered during the proof that were not in the original
specification's list of things to worry about.

\textbf{Sortedness is load-bearing.} The prose specification states
\texttt{BookUncrossed} as the key invariant. The Lean proof required
\texttt{AllInv}---uncrossed plus sorted on both sides---because the
induction step on the matching loop needed to know that the contra side
was monotone in price to establish that the residual order's price was
separated from the contra head. Sortedness is not a cosmetic property.

\textbf{Well-formedness constraints are discovered, not guessed.} The
proof required two fragments of the specification's well-formedness
rules as explicit preconditions: post-only orders have a defined price,
and stop orders are not post-only. Neither was included in the initial
proof attempt. The proof failed at specific branches, and the failures
pointed precisely to these missing hypotheses. The proof told us which
constraints were load-bearing. The precise Lean predicate and its
derivation from the specification's well-formedness constraints are
given in \S\ref{sec:wf}.


% ============================================================================
\section{From Specification to Implementation}
\label{sec:synthesis}
% ============================================================================

The fixed specification---prose amended by the two bug fixes described
in \S\ref{sec:findings}---is designed to serve as the primary input for
a systems-language implementation. The specification's phase structure,
well-formedness constraints, and invariant definitions provide a precise
blueprint: each processing phase maps directly to an implementation
function, each well-formedness constraint maps to an admission check,
and each invariant maps to a runtime assertion.

For the Lean reference implementation, the gap between specification and
implementation is closed by the proof itself: the code that tests run
against is the same code the theorems reference. There is no extraction
layer, no refinement obligation.

For a systems-language implementation such as C++, the gap can be
narrowed by conformance testing against the TLA+ model. In this
approach, TLC generates execution traces---sequences of orders and the
model's expected book state after each---and a test harness replays
these traces against the implementation, comparing its output
state-for-state against the model. This pipeline does not constitute a
refinement proof, but it provides a level of assurance substantially
above conventional unit testing by covering the same combinatorial
feature interactions that TLC explores.

\subsection{The Conformance Pipeline}
\label{sec:conformance-pipeline}

We implemented this approach for a header-only C++ matching engine
(\texttt{matcher\_stl/src/engine.h}, ${\sim}645$ source lines). The
pipeline has three stages. First, TLC is invoked in simulation mode
against seeded configurations of each of twelve scenarios---empty,
shallow and deep book, iceberg-dominant, mixed time-in-force,
stops-pending, self-trade-prevention, and so on---with a depth bound
that keeps each random walk short enough to terminate. Second, a
converter parses each trace and records, for every action, the
spec's projected observables: top-of-book on each side, total
book-order count, pending-stop count, last trade price, and the
list of trades produced by that transition. Third, the C++
conformance harness replays each action against the engine and
compares the engine's observables field-by-field to the projection
after every step.

A sweep of approximately 232{,}000 traces across the twelve
scenarios yielded roughly $1.5\times 10^{6}$ per-step comparisons.
Divergences point directly at the offending transition with a
minimal repro that is typically fewer than ten actions long.

\subsection{A Bug Conformance Testing Found}
\label{sec:cpp005}

\paragraph{The scenario.} A stop-market order is submitted with
time-in-force \textsc{day}; its stop price is not yet met, so it
is placed dormant in the stop set. Later, an amend re-prices a
resting limit order across the spread; the amend executes a trade
whose price satisfies the dormant stop's trigger condition.

The stop fires. Its type is rewritten from \textsc{stop\_market}
to \textsc{market}; its time-in-force is unchanged, as specified.
It enters the matching pipeline as a market order, fills partially
against the remaining liquidity, and reaches the post-match
\emph{dispose} step with a non-zero remainder.

\paragraph{What the specification requires.} The specification is
unambiguous at this point: $\textsc{Dispose}$ drops any market
remainder regardless of time-in-force, and safety invariant INV-8
\textsc{NoRestingMarkets} forbids a market order from appearing
in either book queue. No market order may survive past the
dispose step.

\paragraph{The defect.} The C++ \texttt{Dispose} function gated
the market-drop clause on the \textsc{gtc} branch only; the
\textsc{day} and \textsc{gtd} branches fell through to the
book-insertion path, pushing the leftover market quantity onto
the book at price zero. From that point every subsequent action
carried an extra phantom order visible through the book's
aggregate count, and later aggressors could trade against the
zero-priced order, spuriously producing additional fills.

The correction lifts the market-drop guard to apply to all three
persistent time-in-force branches. The fix preserves the existing
behaviour for user-submitted market orders---which well-formedness
constrains to \textsc{ioc}/\textsc{fok}---and adds the
spec-required coverage for triggered stops, whose converted form
retains the original time-in-force.

\paragraph{Why conventional testing missed it.} The engine had
passed sixty-five directed unit tests covering limit matching,
iceberg reload, FOK feasibility, stop triggering, and partial
fills in isolation, but none exercised a stop-triggers-on-amend
path whose triggered market had leftover quantity. A second
oracle, a sibling C++ reference implementation shadowing the
engine through random differential testing at the scale of tens
of millions of operations per seed, also did not surface the
defect: its own \texttt{dispose} function contained the identical
missing market guard on the \textsc{day}/\textsc{gtd} branches,
and a differential test cannot see a bug both implementations
share. The formal specification, by contrast, treated
\textsc{day}/\textsc{gtd} uniformly with \textsc{gtc} in the
dispose clause and made a resting market a safety violation.
Comparing the engine's behaviour against the spec's projected
state was the first oracle positioned to distinguish the two.

This is the structural argument for spec-based conformance
testing over inter-implementation differential testing. An
independent C++ reference catches all data-structure bugs that
do not appear in both codebases, but it cannot catch a semantic
error that both codebases inherited from a shared mental
model---only a specification rooted outside of C++ can.


% ============================================================================
\section{Discussion}
\label{sec:lessons}
% ============================================================================

\subsection{Feature Composition Is Where Bugs Hide}
\label{sec:lessons-composition}

Both model-checking findings share a structural pattern: each involved
feature is correctly specified in isolation, but the specification omits
a clause covering their interaction. Stop orders are correctly
described; iceberg reload is correctly described; FIFO insertion is
correctly described. None of these sections contains an error when read
alone. The bug lives in the intersection: what happens to an iceberg
order's timestamp when the mechanism that zeros its visible quantity is
not the fill path but the STP DECREMENT path?

This is not a failure of careful writing. It is a structural limitation
of prose specifications. Prose sections have natural scope boundaries.
Interactions between features that span sections are systematically
harder to see than behaviors within a single section.

A formal model collapses all features into a single operational
semantics and evaluates invariants globally. The intersection cases are
not special---they are just more states.

A more robust specification of iceberg reload would state the reload
condition declaratively as an invariant on resting order state:
``a resting order must not have visible quantity~$= 0$ and remaining
quantity~$> 0$.'' Rather than specifying reload as a step triggered by
the fill path, it would be specified as a consequence of this invariant:
any action that could produce this state must be followed by a reload.
This declarative framing would make the omission in DECREMENT
immediately apparent.

\subsection{Model Checking and Theorem Proving Are Complementary}
\label{sec:lessons-complementary}

TLC and Lean surfaced different classes of issue. TLC produced the two
feature-composition defects as concrete multi-step traces within minutes
of starting the exploration. A proof attempt would not have generated
those traces directly; the proof effort is focused on establishing what
holds, not on constructing witnesses for what fails.

Lean's universal quantification, in turn, made the constant-fuel
assumption visible. TLC did not reach configurations where the fuel
bound mattered. Only the universally quantified statement forced the
question.

The cost profiles differ dramatically. TLC returned its first
counterexample in 102 seconds of wall time. The Lean proof accumulated
several weeks of interactive development and produced a 5{,}699-line
artifact. The two are complementary rather than substitutable: TLC's
speed is useful early in specification work; Lean's guarantees scale
beyond any bounded configuration.

\subsection{The Role of AI}
\label{sec:lessons-ai}

The TLA+ model and the Lean~4 proofs were developed interactively with
a large language model acting as a coding assistant. Strategic decisions---the progress-lemma decomposition, the
three-case analysis in the elegant proof, and the side-parameterization
refactor---were human contributions. The bulk of tactic scripting and
operator translation was produced by the model under that guidance.

We mention this because the per-file size and per-lemma count reported
in the paper reflect this mode of work. The methodology described is
independent of whether the implementer is a single engineer, a team, or
an AI-augmented workflow. Neither layer alone would have produced the
artifact: the strategic decompositions guided the model to close proofs
it would not have closed autonomously, and the model produced the
thousands of lines of case-splitting and arithmetic invocations that
the human-directed strategy required.

\subsection{Cost and Value}
\label{sec:lessons-cost}

TLA+ and TLC took approximately two days of work and produced one
concrete invariant counterexample plus one specification gap.
The Lean proof took approximately four weeks and produced a universal
guarantee plus discovered the fuel-bound defect and, during the later
full-invariant extension, an MTL/minimum-quantity transcription defect.
These have different return-on-investment profiles.

For a specification that is still under active development, TLC's rapid
feedback is most valuable: a two-day investment finds concrete
counterexamples before any code is written. For a specification that is
considered stable and whose correctness matters at production scale,
Lean's universal guarantee is the appropriate next step. The two levels
are designed to be used in sequence, not as alternatives.


% ============================================================================
\section{Technical Details}
\label{sec:technical}
% ============================================================================

\subsection{The Specification Source}
\label{sec:spec-source}

The object of study is an internal prose specification,
\texttt{matching-engine-formal-spec.md} version~1.2.0, describing a
generic asset-agnostic order matching engine with price-time priority.
The specification is 979 lines of structured pseudocode covering
primitive domains, order well-formedness constraints, the order book
structure, a core matching algorithm, post-match disposition, five order
types, four STP policies, iceberg reload semantics, stop triggering with
cascading evaluation, cancel, and amend. It was written with care:
well-formedness constraints are enumerated, invariants are stated
formally, and edge cases such as MTL price derivation and FOK dry-run
semantics are treated explicitly.

\subsection{Well-Formedness Constraints}
\label{sec:wf}

The prose specification enumerates 23 well-formedness constraints
(WF-1 through WF-20 plus three a/b suffixed refinements) defining the
valid inputs to the engine. Table~\ref{tab:wf} summarizes them.

\begin{table}[H]
\centering
\caption{Well-formedness constraints (excerpt from the prose
specification). TLA+ \texttt{WellFormed} omits WF-6/7, WF-11/12,
and WF-17; \texttt{MakeOrder} initializes the omitted
\texttt{remainingQty}/\texttt{visibleQty} fields. Lean
\texttt{Order.wellFormed} omits WF-6/7, WF-12, and WF-17; it
includes WF-11 but not WF-12. Lean's \texttt{OrderProcOk}
requires only the two load-bearing fragments (WF-13 and its
corollary).}
\label{tab:wf}
\small
\begin{tabular}{@{}llL{8cm}@{}}
\toprule
ID & Group & Constraint \\
\midrule
WF-1       & Qty     & $\mathit{quantity} > 0$ \\
WF-2/2a/2b & Price   & \textsc{limit}, \textsc{stop\_limit} require
                       $\mathit{price} > 0$; \textsc{mtl} has $\mathit{price}=\bot$ \\
WF-3       & Price   & \textsc{market}, \textsc{stop\_market} have $\mathit{price}=\bot$ \\
WF-4       & Stop    & Stop orders require $\mathit{stopPrice} > 0$ \\
WF-5       & Stop    & Non-stop orders have $\mathit{stopPrice}=\bot$ \\
WF-6/7     & TIF     & \textsc{gtd} requires $\mathit{expireTime} > \mathit{now}$; otherwise $\mathit{expireTime}=\bot$ \\
WF-8/8a    & TIF     & \textsc{market}$\Rightarrow$\textsc{ioc}/\textsc{fok};
                       \textsc{mtl}$\Rightarrow$\textsc{gtc}/\textsc{day} \\
WF-9/10    & Iceberg & $0 < \mathit{displayQty} \leq \mathit{quantity}$; icebergs must be \textsc{limit} \\
WF-11/12   & Iceberg & Initial $\mathit{remainingQty}=\mathit{quantity}$; $\mathit{visibleQty}=\min(\mathit{displayQty},\mathit{quantity})$ \\
WF-13      & PostOnly& Post-only$\Rightarrow$\textsc{limit} \\
WF-14      & PostOnly& Post-only$\not\Rightarrow$\textsc{ioc}/\textsc{fok} \\
WF-15      & PostOnly& \textsc{market}, \textsc{mtl}$\Rightarrow$$\neg$post-only \\
WF-16      & STP     & $\mathit{selfTradeGroup}=\bot \iff \mathit{stpPolicy}=\bot$ \\
WF-17      & STP     & $\mathit{stpPolicy}\in\{\textsc{cancel\_newest},$ $\textsc{cancel\_oldest},$ $\textsc{cancel\_both},$ $\textsc{decrement}\}$ \\
WF-18      & MinQty  & $0 < \mathit{minQty} \leq \mathit{quantity}$ \\
WF-19      & MinQty  & MinQty$\Rightarrow$$\neg$post-only \\
WF-20      & MinQty  & \textsc{fok}$\Rightarrow$$\mathit{minQty}=\bot$ \\
\bottomrule
\end{tabular}
\end{table}

The Lean predicate \texttt{OrderProcOk} used in the main theorem is the
minimum subset of these constraints that the matching loop actually
depends on:

\begin{lstlisting}[language={},basicstyle=\ttfamily\footnotesize]
def OrderProcOk (o : Order) : Prop :=
    (o.postOnly = true -> o.price.isSome = true) /\
    ((o.orderType = .stopLimit \/ o.orderType = .stopMarket) ->
       o.postOnly = false)
\end{lstlisting}

The first conjunct rules out the post-only / no-price configuration
(WF-13: post-only $\Rightarrow$ LIMIT; WF-2: LIMIT $\Rightarrow$ price
defined). The second conjunct---stop orders are not post-only---is a
consequence by contrapositive. Both were discovered during the proof:
earlier attempts omitted them and branches of the matching loop that
case-split on an undefined price failed. The proof told us which WF
clauses were load-bearing.

\subsection{Invariant Suite and Loop Guard}
\label{sec:invariants}

The prose specification states 14 invariants that must hold after every
\texttt{PROCESS} call. Table~\ref{tab:inv-summary} gives the
book-structural core; Table~\ref{tab:invariants} gives the complete
TLC-checked inventory.

\begin{table}[H]
\centering
\caption{Core book invariants.}
\label{tab:inv-summary}
\small
\begin{tabular}{@{}llL{7cm}@{}}
\toprule
ID & Group & Statement \\
\midrule
INV-1   & Structural & Every price level has a non-empty order queue \\
INV-2   & Structural & $\forall$ consecutive bid levels $(L_1,L_2)$: $L_1.\mathit{price} > L_2.\mathit{price}$ \\
INV-3   & Structural & $\forall$ consecutive ask levels $(L_1,L_2)$: $L_1.\mathit{price} < L_2.\mathit{price}$ \\
INV-4   & Structural & $\mathit{bestBid} < \mathit{bestAsk}$ or either side empty (\textbf{uncrossed}) \\
INV-5   & Content    & Every resting order has $\mathit{remainingQty} > 0$ \\
INV-6   & Content    & Every resting order has status $\in\{\textsc{new},\textsc{partially\_filled}\}$ \\
INV-7   & Ordering   & FIFO within each price level: consecutive orders have strictly increasing timestamps \\
INV-8   & Content    & No \textsc{market} orders rest on the book \\
INV-9   & Trade      & Every trade executes at the passive order's price \\
INV-10  & Events     & Events emitted in causal order \\
INV-11  & PostOnly   & No trade where the aggressor was post-only \\
INV-12  & STP        & No trade between orders in the same STP group \\
INV-13  & Content    & No \textsc{mtl} orders rest on the book \\
INV-14  & Content    & No $\mathit{minQty}$ is set on any resting order \\
\bottomrule
\end{tabular}
\end{table}

\begin{table}[H]
\centering
\caption{Invariants checked by TLC and Lean.}
\label{tab:invariants}
\small
\begin{tabular}{@{}llL{4.8cm}ll@{}}
\toprule
ID & TLA+ Operator & Description & TLC & Lean \\
\midrule
INV-1  & \texttt{NoEmptyLevels}    & Every price level has a non-empty order queue & trivial & $\checkmark$ \\
INV-2  & (sortedness)              & Bids strictly descending & by constr. & $\checkmark$ \\
INV-3  & (sortedness)              & Asks strictly ascending & by constr. & $\checkmark$ \\
INV-4  & \texttt{BookUncrossed}    & $\mathit{bestBid} < \mathit{bestAsk}$ or either side empty & $\checkmark$ & $\checkmark$ \\
INV-5  & \texttt{NoGhosts}         & Every resting order has $\mathit{remainingQty} > 0$ & $\checkmark$ & $\checkmark$ \\
INV-6  & \texttt{StatusConsistency}& Resting orders have status \textsc{new} or \textsc{partial} & $\checkmark$ & $\checkmark$ \\
INV-7  & \texttt{FIFOWithinLevel}  & FIFO within level (strictly increasing timestamps) & fixed$^{\star}$ & $\checkmark$ \\
INV-8  & \texttt{NoRestingMarkets} & \textsc{market} orders never appear on the book & $\checkmark$ & $\checkmark$ \\
INV-9  & (construction)            & Trade price equals the passive order's price & by constr. & by constr. \\
INV-11 & \texttt{PostOnlyGuarantee}& No trade where the aggressor was post-only & $\checkmark$ & $\checkmark$ \\
INV-12 & \texttt{STPGuarantee}     & No trade between orders in the same STP group & $\checkmark$ & $\checkmark$ \\
INV-13 & \texttt{NoRestingMTL}     & \textsc{mtl} orders are converted before resting & $\checkmark$ & $\checkmark$ \\
INV-14 & \texttt{NoRestingMinQty}  & $\mathit{minQty}$ is cleared before an order rests & $\checkmark$ & $\checkmark$ \\
\bottomrule
\end{tabular}

\medskip
\raggedright\footnotesize
IDs match the canonical \S13 summary of the prose specification.
TLC column: $\checkmark$ means bounded exhaustive pass across all
configurations in Table~\ref{tab:configs}. Lean column:
$\checkmark$ means mechanically proved for all book sizes---INV-2/3/4
as part of \texttt{AllInv}, the book-state invariants via the capstone
\texttt{process\_preserves\_BookInvariant}, and INV-11/12 via
\texttt{process\_PostOnlyGuarantee} and \texttt{process\_STPGuarantee}.
INV-9 holds by construction; INV-10 (event ordering) is not modeled in
Lean.
$^{\star}$INV-7 was violated under the original specification and
passes after the fix described in \S\ref{sec:bug1}.
\end{table}

A single predicate governs the matching loop's continuation:
\[
\texttt{canMatchPrice}(o, p) \;\equiv\;
\begin{cases}
\top & o.\mathit{type}\in\{\textsc{market},\textsc{stop\_market}\} \\
o.\mathit{price} \geq p & o.\mathit{side}=\textsc{buy},\; o\text{ priced} \\
o.\mathit{price} \leq p & o.\mathit{side}=\textsc{sell},\; o\text{ priced} \\
\top & o.\mathit{type}=\textsc{mtl},\text{ no trades yet} \\
\text{derived price} & o.\mathit{type}=\textsc{mtl},\text{ after first fill}
\end{cases}
\]
This predicate is the loop guard of \texttt{doMatch}. When the loop
terminates in a non-empty contra side, it terminates because
$\neg\texttt{canMatchPrice}(o, \mathit{bestContra})$. This is the seed
of the three-case elegant proof: if matching stopped because prices no
longer cross, then the book cannot be crossed after disposition places
the residual on the same side as the aggressor.

\subsection{TLA+ Model Structure}
\label{sec:model}

The TLA+ specification (\texttt{MatchingEngine.tla}, approximately 950
lines) is a direct translation of the prose specification. The state
consists of eight variables:

\begin{lstlisting}[language=tlaplus]
bidQ           : PRICES -> Seq(Order)   (* per-price bid queues *)
askQ           : PRICES -> Seq(Order)   (* per-price ask queues *)
stops          : Set(Order)             (* dormant stop orders *)
lastTradePrice : PRICES \cup {NULL}     (* last execution price *)
postOnlyOK     : Bool                   (* no post-only aggressor trade *)
stpOK          : Bool                   (* no same-group trade *)
nextId         : Nat                    (* next order ID *)
clock          : Nat                    (* logical clock *)
\end{lstlisting}

The model has four actions: \texttt{SubmitOrder}, \texttt{CancelOrder},
\texttt{AmendOrder}, and \texttt{TimeAdvance}. \texttt{SubmitOrder}
nondeterministically selects all valid order parameters and invokes the
full processing pipeline. TLC's exhaustive exploration therefore
considers every valid combination of order type, side, price, quantity,
TIF, STP policy, iceberg flag, post-only flag, and minimum quantity,
within the configured bounds.

The model is parameterized by \texttt{PRICES}, \texttt{MAX\_QTY},
\texttt{MAX\_ORDERS}, and \texttt{MAX\_CLOCK}. The configurations
explored are shown in Table~\ref{tab:configs}.

\begin{table}[H]
\centering
\caption{Model checking configurations and statistics.}
\label{tab:configs}
\small
\begin{tabular}{@{}lccclrrll@{}}
\toprule
Configuration & Orders & Qty & Prices & Amend & States Gen. & Distinct & Runtime & Result \\
\midrule
Tiny          & 2 & 1 & \{1,2\}   & No  &    2,979,719 &  1,123,333 & 20\,s     & Pass \\
Small         & 2 & 2 & \{1,2\}   & No  &   11,528,343 &  6,037,674 & 1\,m\,42\,s & Pass \\
Medium        & 2 & 2 & \{1,2,3\} & No  &   36,700,016 & 21,261,901 & 5\,m\,34\,s & Pass \\
With Amend    & 2 & 2 & \{1,2\}   & Yes &   25,209,607 &  9,104,902 & 3\,m\,33\,s & Pass \\
3-order       & 3 & 2 & \{1,2\}   & No  & $>$45\,M     & $>$26\,M   & $>$10\,m  & \emph{partial} \\
\bottomrule
\end{tabular}

\medskip
\raggedright\footnotesize
All ``Pass'' results above are post-fix. Prior to fixing the stop-timestamp
bug, the Small configuration produced a counterexample on the
\texttt{FIFOWithinLevel} invariant within the first 11.5 million states. The
3-order run was interrupted before state-space exhaustion; ``partial'' means no
violation was found in the explored states, not clean passage under all
3-order behaviors.
\end{table}

The model does not cover: multiple STP group identifiers (only
\texttt{"G1"} is used); GTD expiry without wall-clock constraints;
event ordering (INV-10); or implementation correctness as opposed to
specification consistency.

\subsection{Lean Proof Architecture}
\label{sec:lean-arch}

The Lean proof artifacts are summarized in Table~\ref{tab:lean-stats}.

\begin{table}[H]
\centering
\caption{Lean~4 proof artifacts (final state).}
\label{tab:lean-stats}
\small
\begin{tabular}{@{}lrrl@{}}
\toprule
File & Lines & Theorems & Proof style \\
\midrule
\texttt{Theorems.lean}         & 3{,}965 & 74 & Constructive, branch-by-branch \\
\texttt{TheoremsFull.lean}     & 1{,}423 & 77 & \S13 book-state suite + trade guarantees \\
\texttt{TheoremsElegant.lean}  &    311  &  6 & Three-case trichotomy \\
\midrule
\textbf{Total}                 & \textbf{5{,}699} & \textbf{157} & Complementary proof artifacts \\
\bottomrule
\end{tabular}
\end{table}

\texttt{Theorems.lean} and \texttt{TheoremsElegant.lean} establish the
\texttt{AllInv} core---uncrossed plus sorted on both sides.
\texttt{TheoremsFull.lean} lifts the remaining \S13 book-state
invariants (no empty levels, no ghosts, status consistency, FIFO, and
no resting market/MTL/MinQty orders) together with the post-only and
STP guarantees on emitted trades, culminating in the capstone
\texttt{process\_preserves\_BookInvariant}. Lifting INV-13 also
surfaced a transcription defect: a well-formed \textsc{mtl} order
carrying a \texttt{minQty} could come to rest mistyped as \textsc{mtl}
at price zero, because the minimum-quantity pre-check short-circuited
the MTL routing phase. The prose specification was correct; the Lean
and TLA+ transcriptions had both diverged from it, and both are now
fixed.

The mutual induction covers three functions---\texttt{processOrder},
\texttt{processCascade}, and \texttt{processTriggeredStops}---which are
mutually recursive because a trade can trigger a dormant stop, which
converts to a market order and re-enters the pipeline, which can in
turn produce more trades and trigger more stops:

\begin{lstlisting}[language={},basicstyle=\ttfamily\footnotesize]
theorem process_all_preserve_AllInv : forall fuel,
  (forall o b, OrderProcOk o -> StopsNoPostOnly b -> AllInv b ->
    AllInv (processOrder fuel o b).book /\
    StopsNoPostOnly (processOrder fuel o b).book) /\
  (forall trades b, StopsNoPostOnly b -> AllInv b ->
    AllInv (processCascade fuel trades b).book /\
    StopsNoPostOnly (processCascade fuel trades b).book) /\
  (forall orders b, (forall o in orders, o.postOnly = false) ->
                    StopsNoPostOnly b -> AllInv b ->
    AllInv (processTriggeredStops fuel orders b).book /\
    StopsNoPostOnly (processTriggeredStops fuel orders b).book)
\end{lstlisting}

Induction proceeds on \texttt{fuel}. The base case (\texttt{fuel = 0})
is trivial because all three functions return the book unchanged. The
step case discharges every phase of \texttt{processOrder}: stop
handling, post-only rejection, FOK pre-check, MinQty pre-check,
market-to-limit routing, normal matching, dispose, and stop cascade.

The subtlest sub-case is MTL Phase~4 case~3: a
\textsc{market-to-limit} order that fills partially, converts to a
\textsc{limit} order, then runs a second matching pass with the derived
limit price. The second \texttt{doMatch} operates on the post-first-match
book with a freshly constructed incoming order, and the non-crossing
precondition for the subsequent dispose must be extracted from a
\texttt{doMatch} call that does not match the form of any existing
helper lemma. The resolution was to parameterize the non-crossing
extraction helpers by the starting timestamp so they apply uniformly to
both the first matching pass (which uses \texttt{b.clock + 1}) and the
MTL second pass (which uses \texttt{mr.clock}).

Every matching lemma has both a buy-side and a sell-side form. The first
proof attempt duplicated every lemma across sides. The fix was to
parameterize each lemma by a \texttt{Side} value and perform
\texttt{cases side} only at the leaves where the concrete direction of
the inequality matters. After side-parameterization, the sell mirror of
the fuel-sufficiency proof compiled on the first try.

\subsection{The Capstone Theorem and Trusted Computing Base}
\label{sec:capstone-tcb}

The book-state guarantees combine into a single capstone theorem. In
full Lean~4 syntax:

\begin{lstlisting}[language={},basicstyle=\ttfamily\footnotesize]
theorem process_preserves_BookInvariant
    (b : BookState) (o : Order)
    (hall : AllInv b) (hpok : OrderProcOk o) (hsnp : StopsNoPostOnly b)
    (hb : BookOk b) (hstops : StopsWF b) (hok : OrderRestOk o) :
    BookInvariant (process b o).book
\end{lstlisting}

\noindent Here \texttt{BookInvariant} bundles uncrossedness, no empty
levels, no ghost orders, status consistency, FIFO within each level, and
the no-resting market/MTL/MinQty conditions (INV-1 and INV-4 through
INV-8, INV-13, INV-14); the sortedness conjuncts (INV-2, INV-3) travel
with \texttt{AllInv}. The two trade-event guarantees are proved
separately and \emph{unconditionally}---an order reaches a matching
phase only after failing the post-only pre-check, so every emitted trade
satisfies them by construction:

\begin{lstlisting}[language={},basicstyle=\ttfamily\footnotesize]
theorem process_PostOnlyGuarantee (b : BookState) (o : Order) :
    PostOnlyGuarantee (process b o).trades      -- INV-11
theorem process_STPGuarantee     (b : BookState) (o : Order) :
    STPGuarantee     (process b o).trades       -- INV-12
\end{lstlisting}

\paragraph{Preconditions.} The hypotheses are logical preconditions of
this preservation step, not domain restrictions smuggled in to make the
theorem true; each is discharged from ordinary inputs.
\texttt{OrderRestOk} and \texttt{OrderProcOk} follow from the
specification's well-formedness predicate
(\texttt{OrderRestOk\_of\_WellFormed}). \texttt{BookOk} and
\texttt{StopsWF} hold of the empty book (\texttt{BookOk\_empty},
\texttt{StopsWF\_empty}) and are preserved by \texttt{process}
(\texttt{process\_preserves\_BookOk}), and \texttt{AllInv} is likewise
inductive; so the whole precondition bundle holds of every book
reachable from an empty start, and the invariant holds along every
execution rather than at isolated states.

\paragraph{Trusted computing base.} The proof rests on an unusually
small base. The project builds with Lean~4 (toolchain \texttt{v4.26.0})
and has \emph{no external dependencies}---in particular it does not use
Mathlib; every module imports only other modules of the project. No
proof file contains \texttt{sorry}, \texttt{admit},
\texttt{native\_decide}, or any custom \texttt{axiom}. Printing the axiom
footprint of the capstone confirms it:

\begin{lstlisting}[language={},basicstyle=\ttfamily\footnotesize]
'process_preserves_BookInvariant' depends on axioms:
    [propext, Classical.choice, Quot.sound]
'process_PostOnlyGuarantee' depends on axioms: [propext, Quot.sound]
'process_STPGuarantee'      depends on axioms: [propext, Quot.sound]
\end{lstlisting}

\noindent These are the three standard axioms of Lean's classical core;
the trade guarantees do not even require choice. Every accepted term is
checked by the Lean kernel, and no external decision procedure or SMT
solver participates in acceptance. The AI-assisted development described
in \S\ref{sec:lessons-ai} therefore affects only the \emph{cost} of
producing the proof, not the basis on which it is trusted: a proof the
kernel rejects does not build.

\subsection{Traceability}
\label{sec:traceability}

Table~\ref{tab:trace} makes the chain from specification clause to TLA+
artifact to Lean artifact explicit for the clauses exercised in the main
theorem.

\begin{table}[H]
\centering
\caption{Spec $\to$ TLA+ $\to$ Lean traceability for the clauses
touched by the main theorem. The Lean column refers to
\texttt{Theorems.lean}, \texttt{TheoremsFull.lean}, or
\texttt{Invariants.lean}.}
\label{tab:trace}
\small
\begin{tabular}{@{}lL{3cm}L{3.2cm}L{4.3cm}@{}}
\toprule
Spec clause & TLA+ artefact & Lean artefact & Role \\
\midrule
WF-13       & \texttt{WellFormed} guard & \texttt{OrderProcOk} field 1 & Main-theorem precondition \\
WF-13 (derived) & \texttt{WellFormed} guard & \texttt{OrderProcOk} field 2 & Main-theorem precondition \\
INV-2/3     & not separately checked$^\dagger$ & \texttt{bidsSortedDescB}, \texttt{asksSortedAscB} & Discovered: needed for induction \\
INV-4       & \texttt{BookUncrossed}    & \texttt{BookUncrossed}, part of \texttt{AllInv} & Main-theorem conclusion \\
INV-5       & \texttt{NoGhosts}         & \texttt{FullBookInv} & Proved preserved \\
INV-6       & \texttt{StatusConsistency}& \texttt{FullBookInv} & Proved preserved \\
INV-7       & \texttt{FIFOWithinLevel}  & \texttt{FullBookInv} & Proved preserved \\
INV-8       & \texttt{NoRestingMarkets} & \texttt{FullBookInv} & Proved preserved \\
INV-11      & \texttt{PostOnlyGuarantee}& \texttt{PostOnlyGuarantee} & Proved (no precond.) \\
INV-12      & \texttt{STPGuarantee}     & \texttt{STPGuarantee} & Proved (no precond.) \\
INV-13      & \texttt{NoRestingMTL}     & \texttt{FullBookInv} & Proved preserved \\
INV-14      & \texttt{NoRestingMinQty}  & \texttt{FullBookInv} & Proved preserved \\
\texttt{canMatch}
            & \texttt{CanMatchPrice}    & \texttt{canMatchPrice}   & Matching loop guard \\
PROCESS \S12& \texttt{ProcessOrder} (11 phases) & \texttt{processOrder} (11 phases) & Mutual induction target \\
\bottomrule
\end{tabular}

\medskip
\raggedright\footnotesize
$^\dagger$Structural sortedness is not stated as a separate TLC
invariant because TLA+ sequences are ordered by construction from a
sorted key. In Lean the sortedness is explicit because book levels are
a \texttt{List PriceLevel} and the induction needs to know both sides
are monotone. This mismatch was discovered \emph{during} the proof:
\texttt{BookUncrossed} alone is not inductive, which forced the
addition of \texttt{AllInv = BookUncrossed $\land$ bidsSortedDescB
$\land$ asksSortedAscB} (with the Boolean predicates applied to the
bid and ask level lists and required to equal \texttt{true}).
\end{table}


% ============================================================================
\section{Related Work}
\label{sec:related}
% ============================================================================

This work sits at the intersection of several literatures: verified
auction and matching algorithms, industrial model checking, mechanized
systems verification, model-based conformance testing, and
machine-checked termination. We survey each and locate our contribution
relative to the closest prior art.

\paragraph{Verified matching engines and auctions.}
The closest prior art is a line of Coq formalizations of exchange
matching by Sarswat, Singh, Natarajan, and Garg. They verify
\emph{double-auction} matching: an early result proves fairness,
uniformity, and individual rationality for single-round call
auctions~\cite{Sarswat2020}, extended to multi-unit trade requests with
uniqueness theorems that justify extracted regulatory
checkers~\cite{Natarajan2021,Garg2025}, and most recently to an
$O(n\log n)$ \emph{continuous} double-auction matcher with a matching
lower bound~\cite{Garg2024}. In industry, Imandra has been used to
verify the matching logic and FIX connectivity of real exchanges and
dark pools~\cite{Passmore2017,Imandra2020,Passmore2021}. Our work is
complementary on three axes. First, \emph{object}: we prove invariant
preservation across a stateful order book over time, rather than the
input/output correctness of one matching round, so time priority
(FIFO), resting, and re-entry become first-class. Second,
\emph{features}: our engine carries the interacting order-type suite
that makes production matchers hard---icebergs, four STP policies,
post-only, FOK, minimum quantity, and stop triggering---and the defects
we report live precisely in their intersections. Third, \emph{method}:
we cross-check a bounded TLA+ model, an unbounded Lean proof, and a C++
implementation against one another, rather than extracting a single
verified reference. Published accounts combining all of these---a
continuous, feature-rich, price-time-priority engine with a mechanized
arbitrary-size proof \emph{and} implementation conformance---remain, to
our knowledge, rare. Adjacent threads---market-design analysis of
frequent batch auctions~\cite{Budish2015} and smart-contract
verification for decentralized exchanges~\cite{Bhargavan2016,%
Grishchenko2018}---are orthogonal to the operational correctness of a
central-limit-order-book engine.

\paragraph{TLA+ and model checking in industry.}
TLA+ has been applied at scale at Amazon Web
Services~\cite{Newcombe2015,Newcombe2014}, at Microsoft (Azure Cosmos
DB)~\cite{Hackett2023}, and at MongoDB~\cite{Davis2020}, with a recent
survey cataloguing a decade of industrial practice~\cite{Bogli2024}. We
follow the Newcombe et~al.\ template---use TLC to explore a
specification's state space and find interaction bugs before
coding---but differ in domain (a financial matching engine rather than a
distributed storage protocol) and in the follow-through to a
machine-checked proof in Lean~4. MongoDB's ``eXtreme
Modelling''~\cite{Davis2020} is the closest methodological analogue to
our conformance stage: it generates tests from a TLA+ model and checks
an implementation's traces against it. We adopt the same spec-as-oracle
idea and combine it with random differential testing.

\paragraph{Model-based and conformance testing.}
Testing an implementation against a formal model has a long theory, from
Tretmans's \textsc{ioco} conformance relation on labelled transition
systems~\cite{Tretmans2008} to the model-based-testing taxonomy and
tooling of Utting, Pretschner, and Legeard~\cite{Utting2012,Utting2007}.
Our conformance harness is a concrete offline-generation,
online-execution instance: TLC counterexample and simulation traces
become the tests, and the model-checked TLA+ specification is the
oracle. We trade \textsc{ioco}'s generality for a mechanized reference
model and a domain-specific projection of observables.

\paragraph{Mechanized proof of systems software.}
The largest machine-checked software artifacts---seL4~\cite{Klein2009},
CompCert~\cite{Leroy2009}, CertiKOS~\cite{Gu2016}, FSCQ~\cite{FSCQ2015},
and CakeML~\cite{CakeML2014}---prove functional correctness or
refinement and consumed person-decades of effort.
IronFleet~\cite{IronFleet2015} is especially relevant: it pairs
TLA-style state-machine refinement with Hoare-logic verification, much
as we pair a TLA+ action model with a Lean invariant proof---though we
keep the levels as independent cross-checks over a separately written
C++ engine rather than a single refinement stack.
Verdi~\cite{Verdi2015} and Iris~\cite{Iris2015} represent, respectively,
verified distributed implementations and a general invariant-based
concurrent program logic. The present work is narrower in ambition and
cost (weeks, not years): we prove preservation of the full \S13
book-state invariant suite for the Lean reference implementation, and
its complexity lives in feature composition rather than in concurrency
or low-level memory semantics.

\paragraph{Termination and well-founded recursion.}
The device in \S\ref{sec:lean-fuel-fix}---replacing an arbitrary fuel
bound with a state-derived measure and a progress lemma---is a standard
pattern for defining general recursive functions and proving them total,
studied extensively for HOL provers~\cite{Krauss2009,Bulwahn2007,%
Krauss2011,Manolios2006} and surveyed across the partiality/recursion
design space---domain predicates, fuel and step-indexing, sized types,
well-founded recursion~\cite{Bove2016,Abel2013,Leino2013}. Automated
methods search for lexicographic or size-based orders~\cite{Bulwahn2007};
our measure is instead a single problem-specific function of book state
whose decrease is also the load-bearing soundness argument for the
matching loop---so proving termination and proving the fuel bound
sufficient are the same step, pushed through a per-branch progress
argument over 14 interacting branches.

\paragraph{Differential and property-based testing.}
Differential testing---comparing independent implementations on random
inputs~\cite{McKeeman1998}---underlies our shadow test, which
cross-checks the engine against a vector-based reference. The technique
is well established for compilers, from Csmith~\cite{Csmith2011} to
equivalence modulo inputs~\cite{EMI2014}. Property-based
testing~\cite{Claessen2000} and its foundational verified variant
QuickChick~\cite{QuickChick2015} connect random testing to stated
properties. The lesson of our C++ finding (\S\ref{sec:cpp005}) sharpens
this literature: a differential test between two implementations cannot
catch a semantic error both inherit from a shared mental model---only a
spec-rooted oracle can.

\paragraph{LLM-assisted formal verification.}
A growing body of work uses large language models to draft proofs and
drive proof search: hammer integration~\cite{Thor2022},
informal-to-formal sketching~\cite{DraftSketch2023}, growing lemma
libraries~\cite{LEGOProver2024}, large-scale synthetic Lean
data~\cite{DeepSeekProver2024}, retrieval-augmented premise
selection~\cite{Yang2023}, whole-proof generation~\cite{First2023}, and
in-context proof agents~\cite{COPRA2024}. These target \emph{autonomous}
tactic search. Our experience is complementary: the model did not close
the proof autonomously; it wrote naturalistic Lean under human strategic
direction. The progress-lemma decomposition, the side-parameterization
refactor, and the three-case trichotomy were human contributions; the
$\sim$5{,}700 lines of case-splitting and \texttt{omega} invocations
were the model's. Neither layer alone would have produced the artifact,
and---per \S\ref{sec:capstone-tcb}---the Lean kernel checks every
accepted term regardless of who or what drafted it.

\paragraph{Spec-to-implementation synthesis.}
Program synthesis from formal specifications is a classical research
area~\cite{Manna1980}. Recent work on LLM-driven
synthesis~\cite{Jain2022} has reached practical levels for small
programs. The approach taken in this work stops short of synthesis or refinement
proof. Instead, it carries the specification forward to a C++
implementation through conformance testing against TLC traces, which is
weaker than a proof but strong enough to expose semantic defects shared
by ordinary unit tests and a sibling C++ reference.


% ============================================================================
\section{Limitations and Reproducibility}
\label{sec:limitations}
% ============================================================================

\subsection{Bounded Exploration, Not Proof (for TLA+)}

TLC performs bounded exhaustive state enumeration. A clean run means no
invariant violation was found within the explored configurations. It
does not mean no violation exists for larger configurations, more
orders, or more price levels.

Specific gaps:
\begin{itemize}[nosep]
\item The 3-order partial run did not complete state-space exploration.
\item No configuration with more than two distinct STP groups was
      checked.
\item GTD time-based expiry was checked only via the simplified
      \texttt{TimeAdvance} action, not against wall-clock constraints.
\item The event-ordering invariant (INV-10) is not formalized in the
      model.
\end{itemize}

\paragraph{Lean closes the unbounded gap for the full \S13 book-state
suite.} The Lean proof establishes
\texttt{process\_preserves\_BookInvariant} for arbitrary book sizes and
arbitrary order counts, under preconditions discharged by
well-formedness and the empty-book base case. This covers the
\texttt{AllInv} core (INV-2, INV-3, INV-4) together with no empty levels
(INV-1), no ghosts (INV-5), status consistency (INV-6), FIFO within
level (INV-7), no resting markets (INV-8), no resting MTL (INV-13), and
no resting MinQty (INV-14); the post-only (INV-11) and STP (INV-12)
guarantees on emitted trades are proved separately and unconditionally.
Only INV-9 (trade price equals the passive price, which holds by
construction) and INV-10 (event ordering, which is not modeled in Lean)
lie outside the Lean development.

\subsection{Specification vs.\ Implementation}

For the Lean reference implementation, the gap is closed. The Lean~4
code that tests run is the same code the theorems reference---there is
no separate model, no extraction layer, and no refinement obligation.
The semantic gap between tested behavior and proved behavior is zero.

For the C++ implementation, a gap remains. The Lean proof does not
transfer automatically to C++, and the TLA+-trace conformance pipeline
compares projected observables rather than proving refinement. The
conformance-testing approach described in \S\ref{sec:synthesis} narrows
this gap and found a real implementation defect, but it does not close
the gap. A refinement proof connecting the Lean reference to the C++
implementation is a natural next step.

\subsection{Reproducibility}

All artifacts---the prose specification, the TLA+ model and
configurations, the Lean~4 project, the C++ implementation and
conformance harness, and the paper source---are available
in a public Git repository at
\url{https://github.com/formally-verified-exchange/verified-matching-engine}.
The top-level \texttt{README.md} lists prerequisites and gives
step-by-step commands for building the Lean proofs and running TLC
against each configuration. The complete artifacts are:

\begin{itemize}[nosep]
\item \texttt{matching-engine-formal-spec.md}---the prose specification.
\item \texttt{matcher\_tla/MatchingEngine.tla}---the TLA+ specification
      ($\sim$950 lines).
\item \texttt{matcher\_tla/MatchingEngine.cfg} and the
      per-configuration files \texttt{MatchingEngine\_*.cfg} for
      Table~\ref{tab:configs}.
\item \texttt{matcher\_tla/REPORT.md}---model checking output and bug
      descriptions.
\item \texttt{matcher\_lean/}---the self-contained Lean~4 project.
      The source files live under \texttt{matcher\_lean/MatchingEngine/};
      the three proof files are \texttt{Theorems.lean} (3{,}965 lines,
      74 theorems), \texttt{TheoremsFull.lean} (1{,}423 lines,
      77 theorems), and \texttt{TheoremsElegant.lean} (311 lines,
      6 theorems). All are built and type-checked by running
      \texttt{lake build} from inside \texttt{matcher\_lean/} and are
      imported from the library root \texttt{matcher\_lean/MatchingEngine.lean}.
\item \texttt{matcher\_stl/}---the C++ implementation, directed tests,
      shadow differential harness, TLA+ trace converter, conformance
      replay harness, and \texttt{BUG\_LOG.md} documenting implementation
      defects found by the testing pipeline.
\end{itemize}

To reproduce the FIFO counterexample:
\begin{enumerate}[nosep]
\item Revert the stop timestamp fix in \texttt{ProcessTriggeredStops}:
      replace \texttt{[ConvertStop(s) EXCEPT !.timestamp = tm]} with
      \texttt{ConvertStop(s)}.
\item Run TLC with the Small configuration
      (\texttt{MAX\_ORDERS=2, MAX\_QTY=2, PRICES=\{1,2\}}, amend
      disabled).
\item TLC will find a violation of \texttt{FIFOWithinLevel} within the
      first several million states.
\end{enumerate}

To reproduce the fuel-unsoundness issue (\S\ref{sec:lean-fuel}):
\begin{enumerate}[nosep]
\item Revert \texttt{computeMatchFuel} in
      \texttt{matcher\_lean/MatchingEngine/Process.lean} to the constant
      \texttt{defaultFuel~:=~100}.
\item Re-run \texttt{lake build}. The proof of
      \texttt{doMatch\_preserves\_AllInv} will fail: the measure
      argument no longer dominates the recursive call bound, and any
      book with more than $100$ contra-side orders produces a concrete
      cross after one application of \texttt{process}.
\end{enumerate}


% ============================================================================
\section{Conclusion}
\label{sec:conclusion}
% ============================================================================

We formalized a feature-rich matching-engine specification in TLA+ and
ran bounded exhaustive exploration across several configurations totaling
over 36 million distinct states. TLC identified one specification
defect---a FIFO violation caused by a triggered stop order retaining its
original submission timestamp---and one specification gap in the
interaction between STP DECREMENT and iceberg visibility. The FIFO
counterexample is a three-order sequence whose minimality is part of the
finding: reading the stop, insertion, and FIFO sections of the prose
specification individually does not expose it, but an operational model
that composes all three does.

We then proved in Lean~4 that the processing pipeline preserves the full
\S13 book-state invariant suite---the \texttt{AllInv} core (uncrossed
plus sorted on both sides) together with no empty levels, no ghosts,
status consistency, FIFO within each level, and no resting
market/MTL/MinQty orders, plus the post-only and self-trade-prevention
guarantees on emitted trades---under preconditions discharged by
well-formedness, for arbitrary book sizes. \texttt{Theorems.lean}
(3{,}965 lines, 74 theorems) proves the \texttt{AllInv} preservation
theorem branch-by-branch; \texttt{TheoremsFull.lean} (1{,}423 lines,
77 theorems) lifts the remaining invariants to the capstone
\texttt{process\_preserves\_BookInvariant}; and
\texttt{TheoremsElegant.lean} (311 lines, 6 theorems) adds a shorter
structural argument for the matching-and-dispose core via a three-case
analysis of the matching termination condition. All three files are
type-checked by every build and contain no open obligations. One issue
surfaced during the proof effort itself: the original arbitrary fuel
bound made the preservation statement false for sufficiently populated
books. Replacing it with a state-derived measure and a progress lemma
restores soundness.

The broader takeaway is that bounded model checking and mechanized
theorem proving play complementary roles on the same specification. TLC
is effective at surfacing concrete feature-composition defects inside
small state bounds at speeds useful during early specification work.
Lean provides guarantees that hold for arbitrary scale and, as a side
effect, exposes pipeline issues that bounded exploration cannot reach.
Neither tool substitutes for the other. The combination---a few days of
model checking followed by weeks of theorem proving---represents a
tractable investment for systems where feature interactions can move
real money in the wrong direction.

All artifacts described in this paper---the prose specification, the
TLA+ model and configurations, the Lean~4 project including all proof
files, and the conformance-tested C++ implementation with its
TLA+-driven trace-replay harness---are publicly available at the
repository linked on the first page.


\section*{Acknowledgments}

The TLA+ model and the Lean~4 proofs were developed interactively with
the assistance of a large language model, using the formal specification
as the primary input. We thank the anonymous reviewers for their
comments.

% ============================================================================
\begin{thebibliography}{99}
% ============================================================================

\bibitem{Newcombe2015}
C.~Newcombe, T.~Rath, F.~Zhang, B.~Munteanu, M.~Brooker, and M.~Deardeuff.
How Amazon Web Services Uses Formal Methods.
\emph{Commun. ACM}, 58(4):66--73, 2015.

\bibitem{Klein2009}
G.~Klein et al.
seL4: Formal Verification of an OS Kernel.
In \emph{Proc. ACM SOSP}, 2009.

\bibitem{Leroy2009}
X.~Leroy.
Formal Verification of a Realistic Compiler.
\emph{Commun. ACM}, 52(7):107--115, 2009.

\bibitem{Gu2016}
R.~Gu et al.
CertiKOS: An Extensible Architecture for Building Certified
Concurrent OS Kernels.
In \emph{Proc. USENIX OSDI}, 2016.

\bibitem{Budish2015}
E.~Budish, P.~Cramton, and J.~Shim.
The High-Frequency Trading Arms Race: Frequent Batch Auctions as a
Market Design Response.
\emph{Quarterly Journal of Economics}, 130(4):1547--1621, 2015.

\bibitem{Bhargavan2016}
K.~Bhargavan et al.
Formal Verification of Smart Contracts: Short Paper.
In \emph{Proc. PLAS}, 2016.

\bibitem{Grishchenko2018}
I.~Grishchenko, M.~Maffei, and C.~Schneidewind.
A Semantic Framework for the Security Analysis of Ethereum Smart
Contracts.
In \emph{Proc. POST}, 2018.

\bibitem{Claessen2000}
K.~Claessen and J.~Hughes.
QuickCheck: A Lightweight Tool for Random Testing of Haskell Programs.
In \emph{Proc. ACM ICFP}, 2000.

\bibitem{Manolios2006}
P.~Manolios and D.~Vroon.
Termination Analysis with Calling Context Graphs.
In \emph{Proc. CAV}, 2006.

\bibitem{Leino2013}
K.~R.~M. Leino.
Developing Verified Programs with Dafny.
In \emph{Proc. ICSE}, 2013.

\bibitem{Yang2023}
K.~Yang et al.
LeanDojo: Theorem Proving with Retrieval-Augmented Language Models.
In \emph{Proc. NeurIPS}, 2023.

\bibitem{First2023}
E.~First, M.~Rabe, T.~Ringer, and Y.~Brun.
Baldur: Whole-Proof Generation and Repair with Large Language Models.
In \emph{Proc. ESEC/FSE}, 2023.

\bibitem{Manna1980}
Z.~Manna and R.~Waldinger.
A Deductive Approach to Program Synthesis.
\emph{ACM TOPLAS}, 2(1):90--121, 1980.

\bibitem{Jain2022}
N.~Jain et al.
Jigsaw: Large Language Models Meet Program Synthesis.
In \emph{Proc. ICSE}, 2022.

\bibitem{Sarswat2020}
S.~Sarswat and A.~K. Singh.
Formally Verified Trades in Financial Markets.
In \emph{Proc. ICFEM}, LNCS 12531, pages 217--232, 2020.

\bibitem{Natarajan2021}
R.~Natarajan, S.~Sarswat, and A.~K. Singh.
Verified Double Sided Auctions for Financial Markets.
In \emph{Proc. ITP}, LIPIcs 193, pages 28:1--28:18, 2021.

\bibitem{Garg2024}
M.~Garg and S.~Sarswat.
Efficient and Verified Continuous Double Auctions.
In \emph{LPAR-25 (Complementary Volume)}, 2024. Also arXiv:2412.08624.

\bibitem{Garg2025}
M.~Garg, N.~Raja, S.~Sarswat, and A.~K. Singh.
Double Auctions: Formalization and Automated Checkers.
\emph{Journal of Automated Reasoning}, 69:17, 2025.

\bibitem{Passmore2017}
G.~O. Passmore and D.~Ignatovich.
Formal Verification of Financial Algorithms.
In \emph{Proc. CADE}, LNCS 10395, pages 26--41, 2017.

\bibitem{Imandra2020}
G.~Passmore, S.~Cruanes, D.~Ignatovich, et al.
The Imandra Automated Reasoning System (System Description).
In \emph{Proc. IJCAR}, LNCS 12167, pages 464--471, 2020.

\bibitem{Passmore2021}
G.~O. Passmore.
Some Lessons Learned in the Industrialization of Formal Methods for
Financial Algorithms.
In \emph{Proc. FM}, LNCS 13047, pages 717--721, 2021.

\bibitem{Newcombe2014}
C.~Newcombe.
Why Amazon Chose TLA+.
In \emph{Proc. ABZ}, LNCS 8477, pages 25--39, 2014.

\bibitem{Davis2020}
A.~J.~J. Davis, M.~Hirschhorn, and J.~Schvimer.
eXtreme Modelling in Practice.
\emph{Proc. VLDB Endowment}, 13(9):1346--1358, 2020.

\bibitem{Hackett2023}
F.~Hackett, J.~Rowe, and M.~A. Kuppe.
Understanding Inconsistency in Azure Cosmos DB with TLA+.
In \emph{Proc. ICSE-SEIP}, pages 1--12, 2023.

\bibitem{Bogli2024}
R.~B\"ogli, L.~Lerena, C.~Tsigkanos, and T.~Kehrer.
A Systematic Literature Review on a Decade of Industrial TLA+ Practice.
In \emph{Proc. iFM}, LNCS 15234, pages 24--34, 2024.

\bibitem{Tretmans2008}
J.~Tretmans.
Model Based Testing with Labelled Transition Systems.
In \emph{Formal Methods and Testing}, LNCS 4949, pages 1--38, 2008.

\bibitem{Utting2012}
M.~Utting, A.~Pretschner, and B.~Legeard.
A Taxonomy of Model-Based Testing Approaches.
\emph{Software Testing, Verification and Reliability}, 22(5):297--312, 2012.

\bibitem{Utting2007}
M.~Utting and B.~Legeard.
\emph{Practical Model-Based Testing: A Tools Approach}.
Morgan Kaufmann, 2007.

\bibitem{Krauss2009}
A.~Krauss.
Automating Recursive Definitions and Termination Proofs in Higher-Order
Logic.
PhD thesis, Technische Universit\"at M\"unchen, 2009.

\bibitem{Bulwahn2007}
L.~Bulwahn, A.~Krauss, and T.~Nipkow.
Finding Lexicographic Orders for Termination Proofs in Isabelle/HOL.
In \emph{Proc. TPHOLs}, LNCS 4732, pages 38--53, 2007.

\bibitem{Krauss2011}
A.~Krauss, C.~Sternagel, R.~Thiemann, C.~Fuhs, and J.~Giesl.
Termination of Isabelle Functions via Termination of Rewriting.
In \emph{Proc. ITP}, LNCS 6898, pages 152--167, 2011.

\bibitem{Bove2016}
A.~Bove, A.~Krauss, and M.~Sozeau.
Partiality and Recursion in Interactive Theorem Provers---An Overview.
\emph{Mathematical Structures in Computer Science}, 26(1):38--88, 2016.

\bibitem{Abel2013}
A.~Abel and B.~Pientka.
Wellfounded Recursion with Copatterns: A Unified Approach to Termination
and Productivity.
In \emph{Proc. ACM ICFP}, pages 185--196, 2013.

\bibitem{IronFleet2015}
C.~Hawblitzel, J.~Howell, M.~Kapritsos, J.~R. Lorch, B.~Parno,
M.~L. Roberts, S.~Setty, and B.~Zill.
IronFleet: Proving Practical Distributed Systems Correct.
In \emph{Proc. ACM SOSP}, pages 1--17, 2015.

\bibitem{Verdi2015}
J.~R. Wilcox, D.~Woos, P.~Panchekha, Z.~Tatlock, X.~Wang, M.~D. Ernst,
and T.~Anderson.
Verdi: A Framework for Implementing and Formally Verifying Distributed
Systems.
In \emph{Proc. ACM PLDI}, pages 357--368, 2015.

\bibitem{FSCQ2015}
H.~Chen, D.~Ziegler, T.~Chajed, A.~Chlipala, M.~F. Kaashoek, and
N.~Zeldovich.
Using Crash Hoare Logic for Certifying the FSCQ File System.
In \emph{Proc. ACM SOSP}, pages 18--37, 2015.

\bibitem{Iris2015}
R.~Jung, D.~Swasey, F.~Sieczkowski, K.~Svendsen, A.~Turon, L.~Birkedal,
and D.~Dreyer.
Iris: Monoids and Invariants as an Orthogonal Basis for Concurrent
Reasoning.
In \emph{Proc. ACM POPL}, pages 637--650, 2015.

\bibitem{CakeML2014}
R.~Kumar, M.~O. Myreen, M.~Norrish, and S.~Owens.
CakeML: A Verified Implementation of ML.
In \emph{Proc. ACM POPL}, pages 179--191, 2014.

\bibitem{McKeeman1998}
W.~M. McKeeman.
Differential Testing for Software.
\emph{Digital Technical Journal}, 10(1):100--107, 1998.

\bibitem{Csmith2011}
X.~Yang, Y.~Chen, E.~Eide, and J.~Regehr.
Finding and Understanding Bugs in C Compilers.
In \emph{Proc. ACM PLDI}, pages 283--294, 2011.

\bibitem{EMI2014}
V.~Le, M.~Afshari, and Z.~Su.
Compiler Validation via Equivalence Modulo Inputs.
In \emph{Proc. ACM PLDI}, pages 216--226, 2014.

\bibitem{QuickChick2015}
Z.~Paraskevopoulou, C.~Hri\c{t}cu, M.~D\'en\`es, L.~Lampropoulos, and
B.~C. Pierce.
Foundational Property-Based Testing.
In \emph{Proc. ITP}, LNCS 9236, pages 325--343, 2015.

\bibitem{Thor2022}
A.~Q. Jiang, W.~Li, S.~Tworkowski, et al.
Thor: Wielding Hammers to Integrate Language Models and Automated
Theorem Provers.
In \emph{Proc. NeurIPS}, 2022.

\bibitem{DraftSketch2023}
A.~Q. Jiang, S.~Welleck, J.~P. Zhou, et al.
Draft, Sketch, and Prove: Guiding Formal Theorem Provers with Informal
Proofs.
In \emph{Proc. ICLR}, 2023.

\bibitem{LEGOProver2024}
H.~Wang, H.~Xin, C.~Zheng, et al.
LEGO-Prover: Neural Theorem Proving with Growing Libraries.
In \emph{Proc. ICLR}, 2024.

\bibitem{DeepSeekProver2024}
H.~Xin, D.~Guo, Z.~Shao, et al.
DeepSeek-Prover: Advancing Theorem Proving in LLMs through Large-Scale
Synthetic Data.
arXiv:2405.14333, 2024.

\bibitem{COPRA2024}
A.~Thakur, G.~Tsoukalas, Y.~Wen, J.~Xin, and S.~Chaudhuri.
An In-Context Learning Agent for Formal Theorem-Proving.
In \emph{Proc. COLM}, 2024.

\end{thebibliography}

\end{document}
